\begin{pastexam}[label={px:ncku-107-cce-05}]{成大電通所 107 學年度 問答題第 1 題}
（問答；12 分）

\textbf{1. Permutation} (12 pts)
\begin{enumerate}[label=\alph*.]
\item (\textbf{3 pts}) Please show the output of the function in serial. Assume the initial state of the
array is \{1 2 3\} and k=0, m =2.
\item (\textbf{7 pts}) Please show the recursive function in order to compute the time complexity of the
function and give the answer of the function. Assume the dimension of the array in $n$ and it takes
O(1) to swap two elements (i.e., swap(a[k], a[i]).
\item (\textbf{2 pts}) Is this program a polynomial time algorithm?
\end{enumerate}
\begin{minted}[fontsize=\small]{cpp}
void Function (char *a, const int k, const int m) {
    if (k == m) { // output permutation
        cout << "[";
        for (int i = 0; i <= m; i++) {
            if (i != m)   cout << a[i] <<", ";
            else          cout << a[i] <<"] "; }
    } else {
        for (i = k; i <= m; i++)) {
            swap(a[k], a[i]);
            Function (a, k+1, m);
            swap(a[k], a[i]); }
    }
}
\end{minted}
\end{pastexam}

\begin{pastsoln}
這份考卷和成大電機所 107 學年度的「資料結構」是同一份（兩系所共用），完整詳解見歷屆試題~\ref{px:ncku-107-ee-ds-05}。

答案：a. \texttt{[1, 2, 3] [1, 3, 2] [2, 1, 3] [2, 3, 1] [3, 2, 1] [3, 1, 2]}。
b. 令 $T(r)$ 是還剩 $r$ 個位置時的時間，$T(1) = c_1 n$（印出 $n$ 個數），$T(r) = r\,T(r-1) + c_2 r$，解得
$T(n) = \Theta(n \cdot n!)$。 c. 不是多項式時間（光是輸出就有 $n!$ 個排列）。
\end{pastsoln}
